\begin{claim}{AC.42}{Proved obstruction to a finite-type-only return bound}
Fixed algebraic coefficient functions evaluated at prime-to-$p$ torsion
points, even when their Frobenius periods are prime to $p$, do not by
themselves give a bounded valuation for the return unit minus one.
More precisely, put $p=2$, let $\omega$ be a primitive cube root of
unity, and set $K=\mathbf Q_2(\omega)$. For $n\ge1$ choose a primitive
$3^n$-th root $\zeta_n$ satisfying
\[
 m_n=3^{n-1},\qquad \zeta_n^{m_n}=\omega.
\]
These roots are prime-to-two Teichmuller points. The fourth-power
Frobenius orbit of $\zeta_n$ has the odd length $m_n$. For the single
fixed degree-one coefficient function $A(X)=1-2X$, its orbit return is
\[
 U_n=\prod_{j=0}^{m_n-1}A(\zeta_n^{4^j})
     =N_{K(\zeta_n)/K}(1-2\zeta_n)
     =1-2^{m_n}\omega.                              \tag{AC.42.1}
\]
In particular
\[
 v_2(A(\zeta_n)-1)=1,\qquad v_2(U_n)=0,
 \qquad v_2(U_n-1)=m_n=3^{n-1}.                    \tag{AC.42.2}
\]
For the full square-Frobenius orbit over $\mathbf Q_2$, the analogous
return is
\[
 N_{\mathbf Q_2(\zeta_n)/\mathbf Q_2}(1-2\zeta_n)
 =\Phi_{3^n}(2)=1+2^{m_n}+2^{2m_n},                 \tag{AC.42.3}
\]
again a unit whose difference from one has valuation $m_n$.
Consequently the annihilator sizes $2^{v_2(U_n-1)}$ grow faster than
$3^{gn}$ for every fixed positive integer $g$.

There is an integral semilinear module realizing these annihilator
sizes. Let $L_n=K(\zeta_n)$, put $M_n=\mathcal O_{L_n}$, and let
$\sigma_n$ be the fourth-power Frobenius of $L_n/K$. Give $M_n$ the
operator
\[
 \mathcal F_n(x)=(1-2\zeta_n)\sigma_n(x).
\]
It normalizes the odd cyclic deck action $\delta_n(x)=\zeta_nx$ by
$\mathcal F_n\delta_n\mathcal F_n^{-1}=\delta_n^4$, and
\[
 \mathcal F_n^{m_n}=U_n,\qquad
 M_n/(1-\mathcal F_n)M_n
       \simeq\mathcal O_K/(2^{m_n}),\qquad
 \exp\bigl((L_n/M_n)^{\mathcal F_n}\bigr)=2^{m_n}.
                                                        \tag{AC.42.4}
\]
This is an algebraic Frobenius-cocycle countermodel to a proposed
intermediate bound. It is not asserted to arise from the Jacobians
in AR.20, or from pure cohomology of Weil weight one. It therefore
does not refute the conditional geometric route in AR.20.
Indeed $|U_n|^2=4^{m_n}+2^{m_n}+1>4^{m_n}$ under a complex
embedding, so the required weight-one equality already fails.
\end{claim}
\begin{proof}
The identity $v_3(4-1)=1$ and the elementary binomial valuation formula
show that $4$ has order $3^{n-1}$ modulo $3^n$. Its powers exhaust
the subgroup of units congruent to one modulo three. Thus the roots
\[
 \{\zeta_n^{4^j}:0\le j<m_n\}
\]
are exactly the roots of $X^{m_n}-\omega$. All prime-to-two roots of
unity belong to unramified extensions of $\mathbf Q_2$, and the
Frobenius over its unramified quadratic extension $K$ is fourth
power. The field degree $[K(\zeta_n):K]$ is therefore $m_n$.
Evaluating the displayed monic polynomial at $X=1/2$ and multiplying
by $2^{m_n}$ proves (AC.42.1). Since $\omega$ and $\zeta_n$ are
two-adic units, (AC.42.2) follows immediately.

The two conjugates over $\mathbf Q_2$ are
$1-2^{m_n}\omega$ and $1-2^{m_n}\omega^2$. Their product is
\[
 (1-2^{m_n}\omega)(1-2^{m_n}\omega^2)
   =1+2^{m_n}+2^{2m_n}.
\]
This is also $\Phi_{3^n}(2)$ by
$\Phi_{3^n}(X)=1+X^{m_n}+X^{2m_n}$. Its difference from one
is $2^{m_n}(1+2^{m_n})$, so the asserted valuation is exact.
Finally $3^{n-1}/n\to\infty$, proving the comparison with $3^{gn}$.

For the integral statement, $M_n$ is free of rank $m_n$ over
$\mathcal O_K$. After extension of scalars to $L_n$, its $m_n$
embeddings diagonalize multiplication, and $\sigma_n$ cyclically
permutes them. Thus $\mathcal F_n$ becomes a weighted cyclic
permutation, with product of its weights $U_n$. Its characteristic
polynomial over $K$ is $T^{m_n}-U_n$, so
\[
                 \det(1-\mathcal F_n)=1-U_n.
\]
Modulo two the operator $\mathcal F_n$ is just the fourth-power
automorphism of $\mathbf F_{4^{m_n}}$. Its fixed space has dimension
one over $\mathbf F_4$. The Smith normal form over the discrete
valuation ring $\mathcal O_K$ therefore has exactly one nonunit
invariant factor, whose valuation equals $v_2(1-U_n)=m_n$.
This proves the cokernel assertion. Applying $1-\mathcal F_n$ to
$0\to M_n\to L_n\to L_n/M_n\to0$, and using its invertibility on
$L_n$, identifies that cokernel with $(L_n/M_n)^{\mathcal F_n}$.
Since $\mathcal O_K$ is unramified quadratic, this finite abelian
group is $(\mathbf Z/2^{m_n}\mathbf Z)^2$.

The coefficient $A$ is invertible at every point under consideration
and has constant reduction one. Thus bounded degree of $A$, bounded
valuation of $A-1$, and periods prime to two all hold in this model.
The growing quantity comes from taking a norm over an increasingly
long orbit. A fixed-locus torsion-distance theorem cannot be applied
to these norm polynomials as though their degrees were bounded.
The exact identities for $1\le n\le6$ are independently checked in
$\mathbf Z[X]/(\Phi_{3^n}(X))$ by
\path{frobenius_norm_obstruction.py}.
\end{proof}
